%=============================================================================!
% 2.10  FROM BALLS TO ATOMS                                                   !
%=============================================================================!
\section{From balls to atoms}
\label{sec:ne-atoms}

Every law of this chapter was found with balls, carts and springs. This
section applies the same laws to atoms and writes them in the form, and in
the units, that molecular dynamics uses.

Atoms are, strictly, quantum objects, and quantum mechanics rather than
Newton's laws describes them exactly. The nuclei, however, are thousands
of times heavier than the electrons, and this book treats them as
classical bodies obeying the second law, while quantum mechanics is used
only to compute the forces on them (Chapter~8). Treating the nuclei
classically is an approximation. It is least accurate for the lightest
nuclei, and for hydrogen the quantum character of the nuclear motion can
change measured properties even at room temperature; Chapter~8 discusses
when the classical picture is adequate.

\subsection*{The equations of motion of $N$ atoms}

Each atom $i$, numbered from 1 to $N$, has a mass $m_i$, a position
$\vb{r}_i$ and a velocity $\vb{v}_i$, and feels a force $\vb{F}_i$ from
the other atoms. We write $\vb{r}^N$ for the whole list of positions,
$\vb{r}_1, \dots, \vb{r}_N$. In the models of this book the force on an
atom depends on where every atom is, but not on how fast the atoms move,
until the thermostats of Chapter~13 add a drag. Before then it is a
function of $\vb{r}^N$ alone, and Chapter~8 explains where it comes
from. The second law for each atom gives
\begin{equation}
   m_i\,\frac{\dd^2\vb{r}_i}{\dd t^2} = \vb{F}_i(\vb{r}^N),
   \qquad i = 1, \dots, N .
   \label{eq:ne-eom-atoms}
\end{equation}
These are $3N$ second-order differential equations, one for each
component of each position, and they are coupled: the equation for one
atom contains the positions of all the others through its force, so none
can be solved alone. Molecular dynamics is the numerical solution of these
equations. By \cref{sec:ne-equations} the state of the system is the list
of all positions together with the list of all velocities, $6N$ numbers
in all, and it determines the motion that follows. The learned
propagators planned for Part~VI, programs trained to predict the state a long
step ahead, take this state as their input for that reason.

\subsection*{Newton's law in eV, \AA, fs and amu}

\Cref{eq:ne-eom-atoms} needs one conversion in the units of this book.
Forces between atoms are measured in electronvolts per \aa ngstr\"om,
\si{\electronvolt\per\angstrom}. The electronvolt is a unit of energy,
which Chapter~3 introduces; for now it is enough that $\SI{1}{\electronvolt}
= \SI{1.602176634e-19}{\newton\metre}$, so that
\si{\electronvolt\per\angstrom} is a unit of force,
\begin{equation*}
   \SI{1}{\electronvolt\per\angstrom}
   = \frac{\SI{1.602176634e-19}{\newton\metre}}{\SI{e-10}{\metre}}
   = \SI{1.602176634e-9}{\newton} ,
\end{equation*}
since $10^{-19}/10^{-10} = 10^{-19-(-10)} = 10^{-9}$.
Masses are measured in atomic mass units, $\SI{1}{\amu} =
\SI{1.66053906892e-27}{\kilogram}$, defined as one twelfth of the mass of
an atom of carbon-12, the commonest form of carbon; tables of atomic
masses use this unit. A force of \SI{1}{\electronvolt\per\angstrom}
acting on a mass of \SI{1}{\amu} therefore produces an acceleration that
we first find in \si{\metre\per\second\squared} and then convert. For the
conversion, multiplying $\SI{1}{\angstrom} = \SI{e-10}{\metre}$ by
$10^{10}$ gives $\SI{1}{\metre} = \SI{e10}{\angstrom}$, and in the same
way $\SI{1}{\second} = \SI{e15}{\femto\second}$, so that $\SI{1}{\second
\squared} = (\SI{e15}{\femto\second})^2 = \SI{e30}{\femto\second\squared}$.
A newton per kilogram is a metre per second squared, by the definition
of the newton, and dividing the powers of ten gives $10^{-9}/10^{-27} =
10^{-9-(-27)} = 10^{18}$, so
\begin{equation*}
   \frac{\SI{1.602176634e-9}{\newton}}{\SI{1.66053906892e-27}{\kilogram}}
   = 0.9648533\times10^{18}\,\si{\metre\per\second\squared}
   = \SI{9.648533e17}{\metre\per\second\squared} ,
\end{equation*}
the last step moving the decimal point. Converting to \aa ngstr\"oms and
femtoseconds,
\begin{equation*}
   9.648533\times10^{17}\times\frac{10^{10}}{10^{30}}\,
   \si{\angstrom\per\femto\second\squared}
   = \SI{9.648533e-3}{\angstrom\per\femto\second\squared} ,
\end{equation*}
since $17 + 10 - 30 = -3$.
This number converts a force in \si{\electronvolt\per\angstrom} divided by
a mass in \si{\amu} into an acceleration in
\si{\angstrom\per\femto\second\squared}.

A lithium atom, of mass \SI{6.94}{\amu}, its standard atomic
weight\cite{Prohaska2022} (the masses of atoms in this book are these
values), under a force of
\SI{1}{\electronvolt\per\angstrom} therefore accelerates at
$9.648533\times10^{-3}/6.94 = \SI{1.390e-3}{\angstrom\per\femto\second
\squared}$. Starting from rest, after \SI{10}{\femto\second} it is moving
at $at = \SI{0.0139}{\angstrom\per\femto\second}$, which, since
$\SI{1}{\angstrom\per\femto\second} = \SI{e5}{\metre\per\second}$
(\cref{sec:mo-units}), is \SI{1390}{\metre\per\second}, about four times
the speed of sound in air, \SI{343}{\metre\per\second},
and it has moved $\tfrac12 at^2 = \SI{0.0695}{\angstrom}$, about
\SI{2}{\percent} of the \SI{3}{\angstrom} between neighbouring atoms in
lithium metal. Atoms are fast, but a few femtoseconds carry them only a
small fraction of the way to their neighbours, which is why molecular
dynamics can follow them in short steps; Chapter~9 shows how the step is
chosen, and a femtosecond is typical. Under the same force a hydrogen
atom, $6.94/1.008 = 6.88$ times lighter, accelerates $6.88$ times faster.

The same numbers show why molecular dynamics leaves gravity out. A
lithium atom has mass $6.94\times\SI{1.66054e-27}{\kilogram} =
\SI{1.152e-26}{\kilogram}$ and so weighs $1.152\times10^{-26}\times
9.80665 = \SI{1.13e-25}{\newton}$. Dividing by
\SI{1.602177e-9}{\newton}, the size of \SI{1}{\electronvolt\per\angstrom},
gives \SI{7e-17}{\electronvolt\per\angstrom}, some $10^{16}$ times
smaller than the force of \SI{1}{\electronvolt\per\angstrom} used above
as an example of a force between atoms.

\subsection*{Newton's law in code}

In \texttt{mdlab} the factor is computed once, from the recommended
values of the physical constants (CODATA), and Newton's law for every
atom at once is then a few lines of code.

\begin{code}
def acceleration(forces: ArrayLike, masses: ArrayLike) -> NDArray:
    forces = np.asarray(forces, dtype=float)
    masses = np.asarray(masses, dtype=float)
    # Each atom's single mass divides all three components of its force.
    return forces / masses[..., np.newaxis] * units.FORCE_TO_ACCEL
\end{code}

The forces arrive as an array of $N$ rows, one per atom, and three
columns, one per component, and the masses as a list of $N$ numbers;
\texttt{np.asarray} turns whatever is passed in, a list for instance, into
arrays of decimal numbers. The dimensions of an array, in this sense, are
the directions along which it is indexed: the forces have two, of
lengths $N$ and 3, and the masses one, of length $N$. When NumPy combines two arrays of different
shapes it lines up their last dimensions, and a dimension of length 1 is
repeated as often as needed to match the other array. Dividing the
$N\times3$ forces by the $N$ masses directly therefore pairs the three
columns of the forces with the $N$ masses. For most $N$ the lengths
differ and NumPy stops with an error, but for $N = 3$ they agree and the
division runs silently with the wrong meaning: the $x$ components of all
three forces are divided by the mass of atom~1, the $y$ components by the
mass of atom~2 and the $z$ components by the mass of atom~3. Writing
\texttt{masses[..., np.newaxis]}, where \texttt{...} stands for every
dimension the array already has and \texttt{np.newaxis} adds a new last
dimension of length 1, turns the list of masses into a single column, an
$N\times1$ array. Its last dimension has length 1, so NumPy
repeats the column across the three components and divides every row of
the forces by the mass in the same row, which is \cref{eq:ne-eom-atoms}
for every atom at once, whatever the value of $N$. Multiplying by \texttt{FORCE\_TO\_ACCEL} gives the accelerations in
\si{\angstrom\per\femto\second\squared}.

\subsection*{Momentum in a simulation}

In an isolated group of atoms every force is internal, and for the
models in this book the forces sum to zero, as the forces on the two carts
did (Chapter~7 shows why this holds even for forces that do not act
between simple pairs). By \cref{eq:ne-momentum} the total momentum
$\vb{P}$ of the atoms is then conserved, and their centre of mass moves at
the constant velocity $\vb{P}/M$.

\begin{practice}
A simulation whose starting velocities have a non-zero total momentum
carries the whole system along at $\vb{P}/M$ for ever, a drift that has
nothing to do with the motion of the atoms relative to one another.
Subtracting $\vb{P}/M$ from every velocity removes it. By
\cref{sec:ne-frames} this is only a change to the frame of reference
moving with the centre of mass, so it alters no acceleration and no
force, and it is the usual first step when starting velocities are
assigned (\cref{sec:md-start}).
\end{practice}

\notebook{02}{Newton's law for a group of atoms, and removing the drift}

\begin{selfcheck}
Under the same force, which accelerates faster, a lithium atom or an
oxygen atom (\SI{15.999}{\amu}), and by what factor?
\end{selfcheck}
